% Generated draft: requires mathematical and source review.
\documentclass[11pt]{article}
\usepackage[margin=25mm]{geometry}
\usepackage{fontspec}
\setmainfont{Latin Modern Roman}
\usepackage{amsmath,amssymb,mathtools,unicode-math}
\setmathfont{Latin Modern Math}
\usepackage{enumitem,needspace,xurl,hyperref}
\setlength{\parindent}{0pt}
\setlength{\parskip}{5pt}
\setlength{\emergencystretch}{3em}
\newcommand{\N}{\mathbb{N}}
\newcommand{\Z}{\mathbb{Z}}
\newcommand{\Q}{\mathbb{Q}}
\newcommand{\R}{\mathbb{R}}
\newcommand{\C}{\mathbb{C}}
\newcommand{\F}{\mathbb{F}}
\newcommand{\A}{\mathbb{A}}
\newcommand{\PP}{\mathbb{P}}
\newcommand{\E}{\mathbb{E}}
\newcommand{\eps}{\varepsilon}
\newcommand{\dd}{\,\mathrm{d}}
\begin{document}
\section*{UnsolvedMath ID: 20001651}
\textbf{Source problem code:} catalogue entry
\subsection{Definitions and mathematical statement}
Let \(F_n\) be the rank-\(n\) free group. The source asks to understand \(\operatorname{Hom}(\operatorname{Out}(F_n),\operatorname{MCG}(S))\), particularly whether an infinite image is possible when \(n\ge4\). It does not fix the surface's genus, punctures, boundary or mapping-class convention. The attached aid's quantitative conclusions concern \(\operatorname{Mod}(S_g)\), the orientation-preserving mapping class group of a closed orientable surface, with \(g\ge2\) for the stated Hurwitz bound.

Write
\[
 \mathcal I_n=\ker\bigl(\operatorname{Out}(F_n)\to\operatorname{GL}(n,\Z)\bigr)
\]
for the outer Torelli subgroup. The aid reports that if \(n\ge5\) and \(\rho:\operatorname{Out}(F_n)\to\operatorname{Mod}(S_g)\) has infinite image, then simultaneously
\[
 (n+1)!\le84(g-1),\qquad
 \rho(\operatorname{Out}(F_n))\text{ is nonamenable},\qquad
 |\rho(\mathcal I_n)|=\infty.
\]
For \(n=4\), it reports the factorial and Torelli conclusions unconditionally, but makes the nonamenability conclusion conditional on the computer-assisted Property \((T)\) theorem stated in Nitsche's \texttt{arXiv:2009.05134}. Its rank-five and higher Property \((T)\) inputs are described as peer reviewed. These conditions are necessary obstructions, not a sufficient condition for an infinite-image homomorphism.

The stronger Torelli statement reported in the aid is: for \(n\ge4\), every finite-index \(\Gamma\le\operatorname{Out}(F_n)\) and homomorphism \(\rho:\Gamma\to\operatorname{Mod}(S_g)\) satisfy
\[
 |\rho(\Gamma\cap\mathcal I_n)|<\infty
 \quad\Longrightarrow\quad |\rho(\Gamma)|<\infty.
\]
The factorial statement pertains to the full source group, not automatically to every finite-index subgroup.

The original remark, attributed only to ``Juan'', warns that the surface and rank need bounds because virtual infinite abelianization can produce an infinite virtually abelian mapping-class image. The aid repairs the quotient construction: if a finite-index subgroup of a finitely generated group \(G\) surjects onto \(\Z\), take its normal finite-index core \(K\). Then \(K\) still has infinite abelianization, and
\[
 G/[K,K]
\]
is an infinite finitely generated virtually abelian quotient, with finite-index abelian subgroup \(K/[K,K]\). The aid reports its embedding into the mapping class group of a sufficiently large, explicitly bounded surface; the actual genus formula is absent from the supplied summary. This is a conditional construction, not a claim that every infinite mapping-class image is virtually abelian. Peer-reviewed Property \((T)\) excludes this mechanism for \(\operatorname{Out}(F_n)\), \(n\ge5\), and the aid excludes it conditionally on the stated preprint for \(n=4\). The three gates and this repair do not classify the full Hom set or settle all surfaces in the original question.
\subsection{Short English statement}
Investigate infinite mapping-class images of \(\operatorname{Out}(F_n)\). For closed targets, the supplied aid gives a genus bound, a Torelli-image obstruction and a nonamenability obstruction, with a separate rank-four preprint caveat, and clarifies the original virtual-abelianization construction without claiming a full classification.
\subsection{Sources}
\begin{itemize}
\item AIM workshop \emph{The geometry of the outer automorphism group of a free group}, October 2010, Section 4, Problem 4.2 and the warning attributed only to ``Juan'': \url{https://aimath.org/WWN/outerauto/outerauto.pdf}.
\item Nitsche, \texttt{arXiv:2009.05134}, the supplied computer-assisted rank-four Property \((T)\) input: \url{https://arxiv.org/abs/2009.05134}. This record treats its use conditionally and supplies no rerun certificate.
\item Local registry \texttt{data/unsolvedmath/problems.json}, permanent ID 20001651, source code \texttt{catalogue entry}; canonical source \texttt{aim-geometric-group-theory-notes.json}, \texttt{notes[159]}.
\item Attached machine-generated attempt 1: partial result, run status partially solved, internal rigor label ``proved'', low novelty confidence; not peer reviewed and requiring independent verification before citation. Other theorem inputs and the claimed explicit larger-surface genus bound have no full references or formulas in the supplied summary.
\end{itemize}
\end{document}
